\section{Discrete Poincare}

The finite reader begins as a chain. A chain is a ledger of neighboring
readings: node after node, difference after difference, with a boundary that
marks where the record has already committed itself. The chain is not yet a
space filled with points. It is a sequence of permitted comparisons. Each
edge asks whether two neighboring nodes differ. Each node carries only what
the grammar has allowed the chain to remember.

Energy is the reading that gathers those neighboring differences. In the
simplest finite form it has the shape
\[
  E(x)^2 = \sum_i (x_{i+1}-x_i)^2 .
\]
The formula is not the foundation. The foundation is the obligation it
records. Every adjacent mismatch spends positive activity. If no adjacent
mismatch is visible, the interior energy reads zero. The chain has then
become flat relative to the differences it is allowed to test.

Flatness alone is too weak. A chain can have zero neighboring differences
while every node has drifted by the same amount. The interior edges cannot
see the common shift, because each edge compares one node with the next. If
all nodes move together, no edge registers a difference. Chapter 02 named
this kind of unread motion a null mode. It is not an error. It is a precise
statement about what the current reader cannot resolve.

This is where the boundary anchor becomes necessary. An anchor pins one
node, or one boundary reading, as already committed. It identifies a place
the chain cannot spend as free motion. If the first node is fixed at zero,
then a common shift is no longer admissible: shifting every node would shift
the pinned node, and the grammar forbids that. The anchor does not add a new
interior force. It forbids the chain from hiding activity in a motion the
interior differences cannot read.

The discrete Poincare reading follows. If the anchored chain has zero
energy, then every neighboring difference is zero. Therefore every node is
equal to its neighbor. The chain is constant. But the anchor says the
constant at the boundary is zero. The only constant compatible with the
anchor is the zero chain. Thus zero energy forces every node to zero.

This is a small theorem with a large grammatical weight. Before the anchor,
zero energy meant only zero activity modulo drift. After the anchor, zero
energy means zero activity itself. The quotient has been resolved by the
boundary commitment. The grammar no longer permits a nonzero chain to pass
as zero simply because its activity is invisible to the edge reader.

The point is not that every physical boundary must be fixed in this literal
way. The point is that a measurement grammar must say which null modes are
being quotiented and which have been forbidden. Sometimes a gauge quotient
identifies the drift as harmless. Sometimes a boundary pin reads the drift
as inadmissible. What cannot be permitted is the ambiguity in which a reader
claims positivity while leaving an unread drift free.

A stiffness example helps because it keeps the chain finite. Imagine three
marks coupled by springs, but read the springs only as a metaphor for
neighboring difference. If all spring differences vanish, the three marks
can still move together unless one mark is pinned. Once it is pinned, every
zero-difference configuration collapses to the pinned value. The energy has
become definite. The concrete operator has acquired the first feature it
needs: it detects zero activity.

This detection is the bridge from burden to scale. In the previous section,
the grammar required one clock to read both answerability and elapsed cost.
Here the chain supplies the finite reader for that requirement. A burden
that spends no activity must be indistinguishable from zero only when the
boundary has removed null drift. Otherwise a costless route could still
carry an unread displacement. The anchored chain forbids that cheap escape.

The discrete Poincare reading therefore grounds the positive invariant in a
finite object. The operator reads neighboring differences. The anchor kills
the nullspace. Zero energy identifies the zero chain. Once this is in
place, energy squared can do more than report activity; it can support a
scale. But a scale is not the same as a square. The grammar must pass
through the square-root door.
